% GENERATED FILE. Edit canonical lessons, then run npm run generate:books.
% canonical-source-hash: fnv1a64:9ca9330e6d4d8118
% canonical-lessons: PT-C502-pescar, PT-C502-pesca, PT-C502-isca, PT-C502-anzol, PT-C502-cacar

\chapter{Pescar e caçar}
\label{ch:pt-pescar-e-cacar}
\hlchaptermodality{502}

\begin{chapteropening}
\textbf{By the end of this chapter:} \emph{I can talk about fishing and hunting.}

Complete the last lesson of chapter 502: I can talk about fishing and hunting.
\end{chapteropening}

\section[pescar]{pescar — to fish; to catch (a fish)}

\label{lesson:PT-C502-pescar}



\begin{quote}
Before the new one: say the Portuguese for watering can, then the Portuguese for flowerpot.
\end{quote}

\subsection*{You'll want to know: pescar}
\textbf{pescar} — \textquotedblleft{}to fish; to catch (a fish)\textquotedblright{}.

\begin{quote}
Ao domingo, o meu tio vai pescar para o rio.
\end{quote}

\begin{quote}
\emph{On Sundays, my uncle goes fishing on the river.}
\end{quote}

\emph{Ir pescar} — to go fishing.

\begin{cousinweb}
\emph{pescar} is inherited, through Vulgar Latin \emph{piscāre}, from Latin \emph{piscārī} \textquotedblleft{}to fish\textquotedblright{} (EtymDB), built on \emph{piscis} \textquotedblleft{}fish\textquotedblright{}.
\end{cousinweb}

\begin{grammarlens}[title={What you've built}]
One more word for telling, explaining and describing at B1.
\end{grammarlens}

\subsection*{Your turn}
\begin{itemize}
\raggedright
  \item \emph{Say it:} \emph{pescar}
  \item \emph{Say it:} \emph{pescar}, once more
  \item \emph{Say it:} a sentence of your own with \emph{pescar}
  \item \emph{Recall:} say the Portuguese for watering can, then the Portuguese for flowerpot, then say \emph{pescar} again
\end{itemize}

\begin{tcolorbox}[breakable,colback=teal!4,colframe=teal!35!black,title={Before you move on}]
What does pescar mean? (\textquotedblleft{}To fish; to catch (a fish)\textquotedblright{}.) Say it once more.
\end{tcolorbox}
\section[a pesca]{a pesca — fishing}

\label{lesson:PT-C502-pesca}



\begin{quote}
Before the new one: say the Portuguese for flowerpot, then the Portuguese for to fish.
\end{quote}

\subsection*{You'll want to know: a pesca}
\textbf{a pesca} — \textquotedblleft{}fishing\textquotedblright{}.

\begin{quote}
A pesca é muito importante nesta vila.
\end{quote}

\begin{quote}
\emph{Fishing is very important in this town.}
\end{quote}

\emph{Barco de pesca} — fishing boat.

\begin{cousinweb}
\emph{pesca} comes from \emph{pescar} \textquotedblleft{}to fish\textquotedblright{} (EtymDB), inherited from Latin \emph{piscārī}.
\end{cousinweb}

\begin{grammarlens}[title={What you've built}]
One more word for telling, explaining and describing at B1.
\end{grammarlens}

\subsection*{Your turn}
\begin{itemize}
\raggedright
  \item \emph{Say it:} \emph{a pesca}
  \item \emph{Say it:} \emph{a pesca}, once more
  \item \emph{Say it:} a sentence of your own with \emph{a pesca}
  \item \emph{Recall:} say the Portuguese for flowerpot, then the Portuguese for to fish, then say \emph{a pesca} again
\end{itemize}

\begin{tcolorbox}[breakable,colback=teal!4,colframe=teal!35!black,title={Before you move on}]
What does a pesca mean? (\textquotedblleft{}Fishing\textquotedblright{}.) Say it once more.
\end{tcolorbox}
\section[a isca]{a isca — bait}

\label{lesson:PT-C502-isca}



\begin{quote}
Before the new one: say the Portuguese for to fish, then the Portuguese for fishing.
\end{quote}

\subsection*{You'll want to know: a isca}
\textbf{a isca} — \textquotedblleft{}bait\textquotedblright{}.

\begin{quote}
O peixe comeu a isca e fugiu.
\end{quote}

\begin{quote}
\emph{The fish ate the bait and got away.}
\end{quote}

Also figurative: \emph{morder a isca} — to take the bait.

\begin{cousinweb}
\emph{isca} is inherited, through Old Portuguese \emph{ysca}, from Latin \emph{ēsca} \textquotedblleft{}food\textquotedblright{} (EtymDB), built on \emph{edere} \textquotedblleft{}to eat\textquotedblright{}.
\end{cousinweb}

\begin{grammarlens}[title={What you've built}]
One more word for telling, explaining and describing at B1.
\end{grammarlens}

\subsection*{Your turn}
\begin{itemize}
\raggedright
  \item \emph{Say it:} \emph{a isca}
  \item \emph{Say it:} \emph{a isca}, once more
  \item \emph{Say it:} a sentence of your own with \emph{a isca}
  \item \emph{Recall:} say the Portuguese for to fish, then the Portuguese for fishing, then say \emph{a isca} again
\end{itemize}

\begin{tcolorbox}[breakable,colback=teal!4,colframe=teal!35!black,title={Before you move on}]
What does a isca mean? (\textquotedblleft{}Bait\textquotedblright{}.) Say it once more.
\end{tcolorbox}
\section[o anzol]{o anzol — fish hook}

\label{lesson:PT-C502-anzol}



\begin{quote}
Before the new one: say the Portuguese for fishing, then the Portuguese for bait.
\end{quote}

\subsection*{You'll want to know: o anzol}
\textbf{o anzol} — \textquotedblleft{}fish hook\textquotedblright{}.

\begin{quote}
O peixe ficou preso no anzol.
\end{quote}

\begin{quote}
\emph{The fish was caught on the hook.}
\end{quote}

Plural \emph{anzóis}.

\begin{cousinweb}
\emph{anzol} is usually traced to a Vulgar Latin diminutive built on Latin \emph{hāmus} \textquotedblleft{}a hook\textquotedblright{} (Lewis \& Short); the exact form is not settled.
\end{cousinweb}

\begin{grammarlens}[title={What you've built}]
One more word for telling, explaining and describing at B1.
\end{grammarlens}

\subsection*{Your turn}
\begin{itemize}
\raggedright
  \item \emph{Say it:} \emph{o anzol}
  \item \emph{Say it:} \emph{o anzol}, once more
  \item \emph{Say it:} a sentence of your own with \emph{o anzol}
  \item \emph{Recall:} say the Portuguese for fishing, then the Portuguese for bait, then say \emph{o anzol} again
\end{itemize}

\begin{tcolorbox}[breakable,colback=teal!4,colframe=teal!35!black,title={Before you move on}]
What does o anzol mean? (\textquotedblleft{}Fish hook\textquotedblright{}.) Say it once more.
\end{tcolorbox}
\section[caçar]{caçar — to hunt}

\label{lesson:PT-C502-cacar}



\begin{quote}
Before the new one: say the Portuguese for bait, then the Portuguese for fish hook.
\end{quote}

\subsection*{You'll want to know: caçar}
\textbf{caçar} — \textquotedblleft{}to hunt\textquotedblright{}.

\begin{quote}
Antigamente, os homens da aldeia caçavam coelhos.
\end{quote}

\begin{quote}
\emph{In the old days, the men of the village hunted rabbits.}
\end{quote}

Hunting is \emph{a caça}; a hunter is \emph{um caçador}.

\begin{cousinweb}
\emph{caçar} is inherited, through Old Portuguese, from Late Latin \emph{captiāre} \textquotedblleft{}to try to catch\textquotedblright{} (EtymDB), built on \emph{capere} \textquotedblleft{}to take\textquotedblright{}.
\end{cousinweb}

\begin{grammarlens}[title={What you've built}]
One more word for telling, explaining and describing at B1.
\end{grammarlens}

\subsection*{Your turn}
\begin{itemize}
\raggedright
  \item \emph{Say it:} \emph{caçar}
  \item \emph{Say it:} \emph{caçar}, once more
  \item \emph{Say it:} a sentence of your own with \emph{caçar}
  \item \emph{Recall:} say the Portuguese for bait, then the Portuguese for fish hook, then say \emph{caçar} again
\end{itemize}

\begin{tcolorbox}[breakable,colback=teal!4,colframe=teal!35!black,title={Before you move on}]
What does caçar mean? (\textquotedblleft{}To hunt\textquotedblright{}.) Say it once more.
\end{tcolorbox}
